\documentclass[11pt,a4paper]{article}

\newcommand{\reportTitle}{Stats -- titre court}
\newcommand{\reportSubtitle}{Période YYYY-MM-DD -- YYYY-MM-DD}

\input{preamble_teal.tex}

\begin{document}

\makeTealTitle
  {Titre principal}
  {Période d'analyse / périmètre}
  {\noindent\begin{infobox}
   \sffamily
   \textbf{Objet} : ...\\[4pt]
   \textbf{Source} : ...
  \end{infobox}}

\section*{Synthèse}

\begin{infobox}[title=Chiffres clés,coltitle=white,colbacktitle=tealPrimary,
                fonttitle=\bfseries\sffamily]
\sffamily
\begin{tabular}{@{}l@{\hspace{2em}}r@{}}
Métrique 1                 & \textbf{N} \\
Métrique 2                 & \textbf{N} \\
\quad sous-catégorie A     & \textbf{N} \\
\quad sous-catégorie B     & \textbf{N} \\
\end{tabular}\\[6pt]
\textbf{Détails qualitatifs} : sortir ici toute ligne dont la valeur droite
serait trop longue pour rentrer dans le tabular ci-dessus.
\end{infobox}

\section*{Détail par catégorie}

\rowcolors{2}{white}{tealPale}
\begin{tabular}{lr}
\rowcolor{tealPrimary}\textbf{\color{white}Catégorie} & \textbf{\color{white}Volume} \\
Catégorie 1   & 10 \\
Catégorie 2   & 7 \\
Catégorie 3   & 3 \\
\end{tabular}

\section*{Observations factuelles}

\begin{infobox}
\sffamily
\begin{enumerate}[leftmargin=*]
\item Observation 1.
\item Observation 2.
\end{enumerate}
\end{infobox}

\end{document}
